% Topology-Aware Sample Sort — presentation slot for CMSC 714 group project.
% Compile via:
%   docker run --rm -v "$PWD:/work" -w /work/src/sample_sort \
%       texlive/texlive:latest pdflatex -interaction=nonstopmode presentation.tex

\documentclass[10pt, aspectratio=169]{beamer}
\usetheme{Madrid}
\usecolortheme{whale}

\usepackage{graphicx}
\usepackage{booktabs}
\usepackage{xcolor}
\usepackage{tikz}

\graphicspath{{../../results/plots/}}

\setbeamertemplate{navigation symbols}{}
\setbeamertemplate{footline}[frame number]

\definecolor{accentgreen}{HTML}{2ECC71}
\definecolor{accentred}{HTML}{E74C3C}
\definecolor{accentblue}{HTML}{3498DB}

\title[Topology-Aware Sample Sort]{Topology-Aware Sample Sort}
\subtitle{CMSC 714 Project --- The bisection-bandwidth axis of the topology study}
\author{Sonok Mahapatra}
\institute[UMD]{Group: J. LaGattuta, C. Dalziel, M. Wijitharathna, S. Mahapatra}
\date{Spring 2026}

\begin{document}

% ----------------------------------------------------------------------
\frame{\titlepage}

% ----------------------------------------------------------------------
\begin{frame}{Project Context}
\textbf{Group study:} how do four MPI workloads perform on five simulated
network topologies? Built under SimGrid 4.1 SMPI; topology defined by XML.

\vspace{0.5em}
\begin{columns}[T]
\column{0.5\linewidth}
\textbf{Workloads}
\begin{itemize}
    \item 3D stencil
    \item 3D FFT
    \item Barnes--Hut N-body
    \item \textcolor{accentblue}{\textbf{Sample sort (mine)}}
\end{itemize}

\column{0.5\linewidth}
\textbf{Topologies}
\begin{itemize}
    \item ring
    \item 3D torus
    \item hypercube
    \item fat tree
    \item dragonfly
\end{itemize}
\end{columns}

\vspace{0.8em}
\textbf{This slot:} the sample-sort row of the matrix.
Sample sort is fundamentally a giant \emph{all-to-all} workload, so it
exercises the network's \alert{bisection-bandwidth axis} more than any
other algorithm in the study.
\end{frame}

% ----------------------------------------------------------------------
\begin{frame}{Sample Sort (PSRS) --- 8 Steps, 1 Bottleneck}
\begin{enumerate}
    \itemsep0.15em
    \item Local sort ($N/p$ keys per rank).
    \item Pick $p-1$ regular samples per rank.
    \item \texttt{MPI\_Allgather} samples (every rank holds $p \cdot (p-1)$ samples).
    \item Sort gathered samples; pick $p-1$ global splitters.
    \item Partition local data by splitters $\to$ \texttt{send\_counts[p]}.
    \item \texttt{MPI\_Alltoall} to exchange the counts.
    \item \alert{\textbf{\texttt{MPI\_Alltoallv}} to redistribute data $\leftarrow$ topology-sensitive}.
    \item Local sort of received bucket.
\end{enumerate}

\vspace{0.5em}
\textbf{Why this is a topology test.} Step 7 is one global all-to-all over
$\sim$8\,MB-per-rank. Bisection bandwidth ($O(1)$ on ring, $O(p)$ on fat
tree) sets a hard ceiling on how fast it can possibly run.
\end{frame}

% ----------------------------------------------------------------------
\begin{frame}{Agnostic Baseline --- Strong Scaling (TOTAL\_N=16M)}
\centering
\includegraphics[width=0.78\linewidth]{agnostic_strong_scaling.png}

\vspace{0.3em}
\small
\textbf{Ring catastrophe} at np=256: $0.083\,$s vs.\ $\sim 0.005\,$s for the
others (\textcolor{accentred}{\textbf{10--20$\times$ gap}}). Predicted by
bisection theory and observed exactly.
\end{frame}

% ----------------------------------------------------------------------
\begin{frame}{Agnostic Baseline --- Weak Scaling (1M keys/rank)}
\centering
\includegraphics[width=0.78\linewidth]{agnostic_weak_scaling.png}

\vspace{0.3em}
\small
Two surprises here. \textbf{(1)} Strong scaling regresses past np=128 across
all topologies --- $O(p^2)$ message-pair latency dominates. \textbf{(2)} At
np=256 \emph{weak}, fat tree is ~$0.149\,$s (almost as slow as ring) ---
\alert{contradicts the simple ``fat tree always wins all-to-all'' intuition}.
\end{frame}

% ----------------------------------------------------------------------
\begin{frame}{Topology-Aware Variants --- One Per Topology}
Step 7 was replaced with a topology-aware redistribution in 5 separate binaries:

\vspace{0.5em}
\begin{table}
\centering
\small
\begin{tabular}{ll}
\toprule
\textbf{Variant} & \textbf{Step 7 algorithm} \\
\midrule
\texttt{ring} & $p-1$ pairwise \texttt{MPI\_Sendrecv}s at distances $1..p-1$ \\
\texttt{torus} & 2-phase row $\to$ column \texttt{Alltoallv} on Cartesian sub-comms \\
\texttt{hypercube} & $\log_2 p$ rounds of recursive halving with per-bucket forwarding \\
\texttt{fat\_tree} & Chunked \texttt{Alltoallv} (4 sub-rounds, smaller in-flight buffers) \\
\texttt{dragonfly} & 2-phase intra-group $\to$ inter-group \texttt{Alltoallv} \\
\bottomrule
\end{tabular}
\end{table}

\vspace{0.6em}
All 5 variants pass \texttt{--verify} on np $\in \{16,32,64,128,256\}$.

\vspace{0.5em}
\textbf{Full grid:} 6 binaries (5 optimized + agnostic) $\times$ 5 topology
XMLs $\times$ 5 np $\times$ 2 modes = \textbf{300 \texttt{smpirun} runs}.
\end{frame}

% ----------------------------------------------------------------------
\begin{frame}{Optimization Results --- Diagonal Speedup, Weak np=256}
\centering
\includegraphics[width=0.72\linewidth]{diagonal_speedup_weak_np256.png}

\vspace{0.2em}
\small
\textbf{The story:} only \textcolor{accentgreen}{fat\_tree-opt} and
\textcolor{accentgreen}{ring-opt} actually help. \textcolor{accentred}{torus-opt}
is \emph{14$\times$ slower} than agnostic on its own topology.
\end{frame}

% ----------------------------------------------------------------------
\begin{frame}{Optimization Results --- Full 6$\times$5 Grid (Weak np=256)}
\centering
\includegraphics[width=0.78\linewidth]{grid_weak_np256.png}

\vspace{0.1em}
\small
Red boxes = matching topology (the diagonal). \textbf{Best per column} isn't
always the matching variant: e.g.\ on \emph{torus}, \texttt{dragonfly-opt}
wins ($0.027\,$s) over \texttt{torus-opt} ($0.342\,$s).
\end{frame}

% ----------------------------------------------------------------------
\begin{frame}{Three Findings}
\textbf{A. Fat-tree chunking solved the np=256 weak surprise.}\\
\quad The mystery slow-down had a measurable cause (in-flight buffer pressure
on the simulated links). Splitting the all-to-all into 4 sub-rounds dropped
it from $0.149\,$s to $0.028\,$s --- \textcolor{accentgreen}{\textbf{5.27$\times$ speedup}}.

\vspace{0.6em}
\textbf{B. Naive 2-phase decomposition is a trap.}\\
\quad The textbook ``row $\to$ column on a 2D grid'' optimization (\texttt{torus-opt})
is \textcolor{accentred}{\textbf{14$\times$ slower}} than agnostic at strong np=256.
SMPI's default \texttt{coll-selector:ompi} is already smart; the explicit
sub-comm setup + double-packing overhead dwarfs the savings.

\vspace{0.6em}
\textbf{C. Algorithmic simplicity beats topology-matching.}\\
\quad Fewer rounds wins, even on the ``wrong'' topology.
\texttt{hypercube-opt} on a \emph{ring} (strong np=256) is $0.055\,$s vs.\
agnostic $0.083\,$s --- 1.5$\times$ faster on a topology it wasn't
designed for.
\end{frame}

% ----------------------------------------------------------------------
\begin{frame}{Conclusion}
\textbf{For the team's broader question} (which topologies suit which algorithms?):
\begin{itemize}
    \item Sample sort cleanly exposes the \alert{bisection-bandwidth axis}: the agnostic ranking (\texttt{fat\_tree} $\approx$ \texttt{hypercube} $\approx$ \texttt{dragonfly} $\ll$ \texttt{torus} $\ll$ \texttt{ring}) is the answer for an algorithm dominated by global all-to-all.
    \item Topology-matched code mostly \emph{preserves} that ranking; it does not unlock dramatically new ones.
\end{itemize}

\vspace{0.7em}
\textbf{The nuanced result of this slice:}
\begin{quote}
\itshape
Topology-aware code only wins when it targets a specific measured
bottleneck --- not when it merely ``matches'' the network structurally.
\end{quote}

\vspace{0.5em}
That is more honest than ``topology-aware is always better'' and more useful
than ``topology-aware never helps.''
\end{frame}

% ----------------------------------------------------------------------
\begin{frame}{Backup --- Per-topology Comparison}
\centering
\includegraphics[width=0.95\linewidth]{per_topology_weak.png}

\vspace{0.2em}
\small Agnostic (black, circles) vs.\ matching-topology optimization (red,
squares) across all np, weak scaling. Fat tree is the only big visible win.
\end{frame}

\end{document}
